% editorial preamble for sources/alco-411-atoms-and-charge-in-type-2025
% Mechanically assembled 2026-10-15 by the source-intake line.
% Wrapper layer follows the LIT precedent sources/alco-351-matching-schur-2024/editorial-preamble.tex;
% author macros are kept VERBATIM from pinned-original-source.tex lines 1-435 (multi-line \def and
% catcode blocks included); document class / metadata / cedram-only commands are not carried over;
% pythonhighlight is replaced by a minimal shim (not in TeX Live).
% NOTE (delivery rule): comments in delivered preambles must stay ASCII-only (CJK in delivered text
% triggers an automatic ctex injection in prepared_delivery, which is fatal on hosts without ctex);
% and no orphan metadata arguments may remain (a dangling \title argument breaks \begin{document}).
\PassOptionsToPackage{dvipsnames,svgnames,table}{xcolor}
\documentclass[11pt]{article}
\usepackage[a4paper,margin=22mm]{geometry}
\usepackage[no-math]{fontspec}
\setmainfont{DejaVu Serif}
\usepackage{amsmath,amssymb,amsthm,mathtools,graphicx,xcolor,hyperref,xurl,xspace,ifthen,enumerate}
\usepackage[shortlabels]{enumitem}
\setlength{\emergencystretch}{4em}
% ---- pythonhighlight shim (publisher ships pythonhighlight.sty; not in TeX Live) ----
\newenvironment{python}{\begin{quote}\ttfamily\small}{\end{quote}}
% ---- theorem environments (shared theo counter, matching the publisher numbering) ----
\newtheorem{theo}{Theorem}[section]
\newtheorem{coro}[theo]{Corollary}
\newtheorem{defi}[theo]{Definition}
\newtheorem{exam}[theo]{Example}
\newtheorem{lemma}[theo]{Lemma}
\newtheorem{prop}[theo]{Proposition}
\newtheorem{rema}[theo]{Remark}
\newtheorem{theorem}[theo]{Theorem}
\newtheorem{conjecture}[theo]{Conjecture}
\newtheorem{claim}[theo]{Claim}   % author-declared after \begin{document} in the source (L438); same counter
\newtheorem*{thm*}{Theorem}
% ---- author macros (verbatim from pinned-original-source.tex) ----
%% The class cedram-ALCO is just a wrapper around amsart.cls (version 2)
%% implementing the layout of the journal, and some additionnal
%% administrative commands. 
%% You can place one option:
%% * "Unicode" if the file is UTF-8 encoded.
%% Here you might want to add some standard packages if those
%% functionnalities are required.


\usepackage{mathtools}
\usepackage{tikz}
\usetikzlibrary{decorations.markings}
\usetikzlibrary{shapes.geometric, arrows}
\usepackage{pgfplots}
\usepackage{subcaption}
%\usepackage{biblatex}
%\pagestyle{empty}
\pgfdeclarelayer{edgelayer}
\pgfdeclarelayer{nodelayer}
\pgfsetlayers{edgelayer,nodelayer,main}

\tikzstyle{none}=[inner sep=0pt]
% Node styles
\tikzstyle{red}=[fill=red, draw=red, shape=circle,minimum size=1mm]
\tikzstyle{blue}=[fill=cyan, draw=cyan, shape=rectangle, minimum size=2mm]
\tikzstyle{sred}=[fill=red!20, draw=red, shape=circle, minimum size=1pt]
\tikzstyle{sblue}=[fill=cyan, draw=cyan, regular polygon, regular polygon sides=3, rotate=180, minimum size=1pt]
\definecolor{mygreen}{RGB}{4,200 ,7}
\tikzstyle{green}=[-, draw=mygreen]
\tikzstyle{black}=[-, fill={rgb,255: red,251; green,255; blue,5}, fill opacity = 0.4]
\pgfplotsset{compat=1.15}
\usepackage{mathrsfs}
\usepackage{cleveref}

\Crefname{claim}{Claim}{Claims}
\Crefname{theo}{Theorem}{Theorems}
\Crefname{prop}{Proposition}{Propositions}
\Crefname{conj}{Conjecture}{Conjectures}
\Crefname{ques}{Question}{Questions}
\Crefname{coro}{Corollary}{Corollaries}
\Crefname{lemm}{Lemma}{Lemmas}
\Crefname{defi}{Definition}{Definitions}
\Crefname{nota}{Notation}{Notations}
\Crefname{rema}{Remark}{Remarks}
\Crefname{exam}{Example}{Examples}
\Crefname{section}{Section}{Sections}
\Crefname{subsection}{Subsection}{Subsections}
\Crefname{figure}{Figure}{Figures}

\newcount\cols
{\catcode`,=\active\catcode`|=\active
 \gdef\Young#1{\hbox{$\vcenter
 {\mathcode`,="8000\mathcode`|="8000
  \def,{\global\advance\cols by 1 &}%
  \def|{\cr
        \multispan{\the\cols}\hrulefill\cr
        &\global\cols=2 }%
  \offinterlineskip\everycr{}\tabskip=0pt
  \dimen0=\ht\strutbox \advance\dimen0 by \dp\strutbox
  \halign
   {\vrule height \ht\strutbox depth \dp\strutbox##
    &&\hbox to \dimen0{\hss$##$\hss}\vrule\cr
    \noalign{\hrule}&\global\cols=2 #1\crcr
    \multispan{\the\cols}\hrulefill\cr%
   }
 }$}}
\gdef\Skew(#1:#2){\hbox{$\vcenter
{\mathcode`,="8000\mathcode`|="8000
  \dimen0=\ht\strutbox \advance\dimen0 by \dp\strutbox
  \def\boxbeg{\vbox
    \bgroup\hrule\kern-0.4pt\hbox to\dimen0\bgroup\strut\vrule\hss$}%
  \def\boxend{$\hss\egroup\hrule\egroup}%
  \def,{\boxend\boxbeg}%
  \def|##1:{\boxend\vrule\egroup\nointerlineskip\kern-0.4pt
    \moveright##1\dimen0\hbox\bgroup\boxbeg}%
  \def\\##1\\##2:{\boxend\vrule\egroup\nointerlineskip\kern-0.4pt
    \kern ##1\dimen0\moveright##2\dimen0\hbox\bgroup\boxbeg}%
  \moveright#1\dimen0\hbox\bgroup\boxbeg#2\boxend\vrule\egroup
 }$}}
}
\newcommand{\halfcirc}[2]{
    \draw [fill=cyan] (#1,#2) circle (2.5pt); 
    \fill[red] (#1,#2) --++ (180:2.5pt) arc (180:0:2.5pt) -- cycle; 
}

\newcommand{\bbQ}{\mathbb{Q}}
\newcommand{\bbZ}{\mathbb{Z}}
\newcommand{\bbN}{\mathbb{N}}
\newcommand{\bbC}{\mathbb{C}}
\newcommand{\bbR}{\mathbb{R}}
\newcommand{\bbT}{\mathbb{T}}

\newcommand{\raw}{\rightarrow}
\newcommand{\ra}{\rightarrow}

\newcommand{\affphi}{\widehat{\phi}}

\newcommand{\affell}{\widehat{\ell}}

\newcommand{\calA}{\mathcal{A}}
\newcommand{\calB}{\mathcal{B}}
\newcommand{\calP}{\mathcal{P}}
\newcommand{\calS}{\mathcal{S}}
\newcommand{\calN}{\mathcal{N}}
\newcommand{\calD}{\mathcal{D}}
\newcommand{\affD}{\widehat{\calD}}


\newcommand{\Gr}{\mathcal{G}r}
\newcommand{\sch}[1]{\overline{\Gr_{#1}}}

\newcommand{\affW}{\widehat{W}}
\newcommand{\affX}{\widehat{X}}
\newcommand{\affPhi}{\widehat{\Phi}}
\newcommand{\affDelta}{\widehat{\Delta}}
\newcommand{\affeps}{\widehat{\epsilon}}
\newcommand{\affe}{\widehat{e}}
\newcommand{\afff}{\widehat{f}}

\newcommand{\eps}{\epsilon}
\newcommand{\op}{\operatorname}

\newcommand{\undc}{\underline{c}}
\newcommand{\extW}{\widehat{W}_{ext}}
\newcommand{\lam}{\lambda}
\newcommand{\lamk}{\lambda-k\varpi_2}
\newcommand{\lama}{\lambda-a\varpi_2}
\newcommand{\lamA}{\lambda-A\varpi_2}
\newcommand{\bPsi}{\overline{\Psi}}

\newcommand{\tightoverset}[2]{%
  \mathop{#2}\limits^{\vbox to -.2ex{\kern-0.75ex\hbox{$#1$}\vss}}}
  

\newcommand{\lce}{\left\lceil}
\newcommand{\lfl}{\left\lfloor}
\newcommand{\frf}[1]{\left\lfloor\frac{#1}{2}\right\rfloor}
\newcommand{\frc}[1]{\left\lceil\frac{#1}{2}\right\rceil}

\newcommand{\rce}{\right\rceil}
\newcommand{\rfl}{\right\rfloor}

\newcommand{\calH}{\mathcal{H}}
\newcommand{\bfH}{\mathbf{H}}
\newcommand{\undH}{\underline{\bfH}}
\newcommand{\bfN}{\mathbf{N}}
\newcommand{\tilN}{\widetilde{\bfN}}
\newcommand{\htil}{\widetilde{h}}
\newcommand{\ttt}{\mathtt{T}}
\newcommand{\ttu}{\mathtt{U}}
\newcommand{\Text}{T^\vee_{\text{ext}}}


\DeclareMathOperator{\Ima}{Im}
\DeclareMathOperator{\Arr}{Arr}
\DeclareMathOperator{\sgn}{sgn}
\DeclareMathOperator{\wt}{wt}
\DeclareMathOperator{\Conv}{Conv}
\DeclareMathOperator{\cha}{ch}
\DeclareMathOperator{\pat}{pat}
\DeclareMathOperator{\at}{at}
\DeclareMathOperator{\str}{str}
\DeclareMathOperator{\HL}{HL}
\DeclareMathOperator{\word}{word}
\DeclareMathOperator{\grdim}{grdim}

\makeatletter
\newcommand{\mylabel}[2]{#2\def\@currentlabel{#2}\label{#1}}
\makeatother

\newcommand{\stn}{\str_2}




%%%%%index making commands


%\usepackage[acronym,symbols,sort=use]{glossaries}
%
%\makeglossaries



% \newglossaryentry{sample}{name={sample},description={an example}}

% The $U(\mathfrak{sp}(2n,\mathbb{C}))$-crystal of Kashiwara--Nakashima tableaux of shape $\lambda$
\newcommand{\symptab}{\ensuremath{KN(\lambda,n)}}

% Bruhat graph for a dominant integral weight $\lambda$
\newcommand{\bruhatgraph}{\ensuremath{\Gamma_\lambda}}

 % Twisted Bruhat graph
\newcommand{\twistedbruhat}{\ensuremath{\Gamma^{m}_{\lambda}}}

  % Inversion set of $y_m$
\newcommand{\inversionsetym}{\ensuremath{N(y_m)}}

% Set of arrows pointing to $\mu$ in $\Gamma^{m}_{\lambda}$
\newcommand{\arr}{\ensuremath{ \Arr_m(\mu,\lam)}}

% Twisted length, i.e. $|\Arr_m(\mu,\lam)|$
\newcommand{\twistedlength}{\ensuremath{\ell_m(\mu,\lam)}}

 % $m$-twisted length
\newcommand{\twistedlengthv}{\ensuremath{\ell_m(v)}}

% Number of non-swappable edges $\mu \rightarrow t_k \mu$, for $\mu < t_k \mu < \lambda$ and $k \leq m$
\newcommand{\nm}{\ensuremath{\calN_m(\mu,\lambda)}}

% $\min(k,\affD_m(\mu,\lam-k\varpi_2))$
\newcommand{\dmk}{\ensuremath{\calD_m(\mu,\lam,k)}}
  
    % $X \oplus \bbZ$, the cocharacter lattice of $\Text$
\newcommand{\affx}{\ensuremath{ \affX}}

   % renormalized Kazhdan--Lusztig polynomials
   % \htil is defined above and is just \tilde{h}
%\newcommand{\htil}{\ensuremath{\htil^\eta_{\mu,\lam}(v)}}

    % The extended torus $T^{\vee} \times \bbC^{*}$
\newcommand{\torusext}{ \ensuremath{\Text}}

    % The subset $X_+\subset X$ of dominant weights, that is, $\lambda \in X: \langle \lambda, \alpha^{\vee}\rangle \geq 0$ for $\alpha^{\vee}$ a positive coroot
\newcommand{\dom}{ \ensuremath{X_+}}

    % The complex reductive group $Sp_4(\bbC)$
\newcommand{\g}{ \ensuremath{G}}

    % The Langlands dual of $G$, i.e.  $G^\vee=SO_5(\bbC)$.
\newcommand{\gdual}{ \ensuremath{G^{\vee}}}

     % A choice of Borel subgroup $B^{\vee}\subset G^{\vee}$
\newcommand{\bdual}{ \ensuremath{B^{\vee}}}

     % Maximal torus in $B^{\vee}$
\newcommand{\tdual}{ \ensuremath{T^{\vee}}}

     % The weight lattice, i.e.  the set of cocharacters $ \operatorname{Hom}(\mathbb{C}^{*}, T^{\vee})$
\newcommand{\X}{ \ensuremath{X}}

      % The coweight lattice, i.e.  the set of characters $\operatorname{Hom}(T^\vee,\mathbb{C}^{*})$ of $T^\vee$
\newcommand{\Xvee}{ \ensuremath{X^\vee}}

       % The root system of type $C_2$
       % \Phi already defined
%\newcommand{\Phi}{\ensuremath{\Phi}}

     % The dual root system of $\Phi$
\newcommand{\Phivee}{ \ensuremath{\Phi^\vee}}

% A fundamental weight in $X$
\newcommand{\fw}{\ensuremath{\varpi_1}}

% A positive root in $\Phi$
\newcommand{\alphaone}{\ensuremath{\alpha_1}}

% A positive coroot in $\Phi^\vee$
\newcommand{\alphavee}{\ensuremath{\alpha^\vee_1}}

% A fundamental coweight in $X^\vee$
\newcommand{\dualfw}{\ensuremath{\varpi^\vee_1}}

% The half sum of the positive roots and coroots, respectively
% \rho is already defined
%\newcommand{\rho}{\ensuremath{\rho}}

% The Weyl group of type $C_2$
\newcommand{\W}{\ensuremath{W}}

% The affine Weyl group of type $C_2$. It is the semidirect product of $W$ and $\bbZ \Phi$
% already defined above
%\newcommand{\affW}{\ensuremath{\affW}}

% The affine Weyl group of type $C_2$. It is the semidirect product of $W$ and $X$
% already defined above, also probably the definition below was incorrect. It should be: \widehat{W}_{ext}
%\newcommand{\extW}{\ensuremath{\affW}}

     % The normal crystal associated to the irreducible representation of $\mathfrak{g}$ of highest weight $\lambda$
\newcommand{\blambda}{ \ensuremath{ \calB(\lambda)}}

% The subset of elements of weight $\mu$ in the crystal $\calB(\lambda)$
\newcommand{\blambdamu}{ \ensuremath{ \calB(\lambda)_{\mu}}}

     % The length of the longest non-swappable staircase over $(\mu,\lam)$ where the label of every edge in $e_i$ is a root in $N(y_m)$
\newcommand{\dm}{ \ensuremath{ \affD_m(\mu,\lam)}}

     % The length of the longest non-swappable staircase over $(\mu,\lam)$
\newcommand{\dminfty}{ \ensuremath{ \affD_{\infty}(\mu,\lam)}}

     % The principal preatom $\calP(\lambda) \subset \calB(\lambda)$
\newcommand{\preatom}{ \ensuremath{ \calP(\lambda)}}

     % The atom $\mathcal{A}(\lambda) \subset \mathcal{P}(\lambda)$
\newcommand{\atom}{ \ensuremath{ \calA(\lambda)}}

     % The atomic number
\newcommand{\atomicnumber}{ \ensuremath{ \at(\ttt)}}

     % The preatomic number, i.e.  the maximum integer with $\ttt\in \Ima(\Phi^{\pat(\ttt)})$
\newcommand{\preatomicnumber}{ \ensuremath{ \pat(\ttt)}}

     % Number of non-swappable edges $\mu\raw t_k\mu$ for $ \mu<t_k\mu \leq \lambda$ and $k \leq m $ 
\newcommand{\nonswap}{ \ensuremath{\calN_m(\mu,\lambda)}}

    % Number of non-swappable edges $\mu\raw t_k\mu$ for $ \mu<t_k\mu \leq \lambda, k \in \mathbb{N}$  
\newcommand{\nonswapi}{ \ensuremath{\calN_{\infty}(\mu,\lambda)}}

% For $i\in \{1,2,21,12\}$, $\affphi_i(\mu,\lambda)$ is the maximum integer $k$ such that $\mu-k\alpha_i\leq \lambda$ 
\newcommand{\affphii}{ \ensuremath{\affphi_i(\mu,\lambda)}}

% The reflection subgroup  $W^{k}:=y_{k}Wy_{k}^{-1}\subset \affW$ for $k\geq 0$ 
\newcommand{\Wk}{ \ensuremath{W^{k}}}

% Reflection corresponding to $M\delta-\alpha_2^\vee$
\newcommand{\vm}{\ensuremath{v_M}}

 % Reflection corresponding to $M\delta-\alpha_{12}^\vee$ 
\newcommand{\qm}{\ensuremath{q_M}}

  % Reflection corresponding to $2M\delta-\alpha_{21}^\vee$
\newcommand{\rsubm}{\ensuremath{r_M}}

  % Bruhat order
\newcommand{\bruhatorder}{\ensuremath{<}}

  % m-twisted Bruhat order
\newcommand{\mtwistedbruhatorder}{\ensuremath{\leq_m}}

  %  $\Gamma_{\lambda}^\infty:=\Gamma^m_\lambda$ for $m\gg 0$ (the twisted graphs $\Gamma^{m}_{\lambda}$ stabilize for large $m$)
\newcommand{\stablebruhatgraph}{\ensuremath{\Gamma_{\lambda}^\infty}}

  % The set of edges in $\Gamma_\lam$
\newcommand{\elambda}{\ensuremath{E(\lambda)}}

  % The set of swappable edges in $\Gamma_\lam$
\newcommand{\es}{\ensuremath{E^S(\lambda)}}

  % The set of non-swappable edges in $\Gamma_\lam$
\newcommand{\en}{\ensuremath{E^N(\lambda)}}

  % The weight map $\wt:\ \calB \longrightarrow X$ in a seminormal crystal $\calB$
  % \wt already defined above
%\newcommand{\wt}{\ensuremath{\wt}}
 
% The crystal operators $e_{i},f_{i}: \calB \longrightarrow \calB \sqcup \left\{ 0\right\}$
  \newcommand{\crystalops}{\ensuremath{e_{i},f_{i}}}

 % $\max\{a \in \mathbb{Z}_{\geq 0} :e_i^a(b)\neq 0\}$ 
\newcommand{\epsi}{\ensuremath{\eps_i(b)}}

 % $\max\{a \in \mathbb{Z}_{\geq 0} :f_i^a(b)\neq 0\}$ 
\newcommand{\phii}{\ensuremath{\phi_i(b)}}

 % A Kashiwara--Nakashima tableau of type $C_2$
\newcommand{\kn}{\ensuremath{\ttt}}

 % The embedding of crystals $\Phi: \calB(\lambda) \rightarrow \calB(\lambda + 2 \varpi_{1})$ 
\newcommand{\preemb}{\ensuremath{\Phi}}
 
 % The preatomic function $Z(\ttt):=\phi_1(\ttt)+\phi_2(\ttt)+\phi_{21}(\ttt)$
\newcommand{\preatomicfn}{\ensuremath{Z(\ttt)}}

 % The weight-preserving map  $\bPsi:\mathcal{P}(\lambda) \hookrightarrow \mathcal{P}(\lambda + \varpi_{2})$
\newcommand{\atomicemb}{\ensuremath{\bPsi}}

%  The total contribution of the swapping functions along the direction $\alpha$, $\alpha\in \Phi_+$
\newcommand{\deltaalpha}{\ensuremath{\Delta^\alpha}}

%  $\ell_m(\mu,\lama)-\calN_m(\mu,\lama)+\calD_m(\mu,\lambda,a)+a+2p$
\newcommand{\sigmam}{\ensuremath{\sigma_m(\ttt)}}

%  $\sigma_m(\ttt)+\langle \lam',\rho^\vee\rangle=-\sigma_m(\ttt)+\langle \lam,\rho^\vee\rangle+3p$
\newcommand{\rem}{\ensuremath{r_m(\ttt)}}


\DeclarePairedDelimiter\abs{\lvert}{\rvert} %Something useful only for this sample's sake: you can erase this line in your file (or find it useful...)
